% Teil II: durchgerechneter Weg.

\subsection*{Gegeben und gesucht}
Abstand zweier Momente maximaler Lautstärke (Schwebungsmaxima):
$T_{\max}=\qty{2}{s}$. Gesucht: Frequenzdifferenz $\Delta f$.

\subsection*{Ansatz}
Die beiden Instrumente spielen nahezu gleiche Frequenzen $f_1\approx f_2$; die
Lautstärke moduliert mit der Differenzfrequenz. Zwei aufeinanderfolgende Maxima
liegen um $1/\Delta f$ auseinander.

\subsection*{Rechnung}
\begin{equation}
  T_{\max}=\frac{1}{\Delta f}
  \;\Rightarrow\;
  \Delta f=\frac{1}{T_{\max}}=\frac{1}{\qty{2}{s}}
  =\boxed{\;\qty{0.5}{Hz}\;}.
\end{equation}

\begin{ergebnis}
Die Frequenzen der beiden Töne unterscheiden sich um $0.5\,\mathrm{Hz}$
(die Instrumente sind um einen halben Schwebungszyklus pro $2\,\mathrm s$
verstimmt).
\end{ergebnis}

\begin{kontrolle}
\textbf{Einheiten:} $Hz=1/s$, also $1/T_{\max}$ in \unit{\hertz}. \checkmark

\textbf{Grössenordnung:} $0.5\,\mathrm{Hz}$ ist eine langsame, gut hörbare
Schwebung --- zwei Maxima pro $2\,\mathrm s$ entspricht gerade dem Gehörten.
\end{kontrolle}

\begin{merksatz}
Schwebungsfrequenz $=\Delta f$; der Abstand zweier Lautstärke-Maxima ist
$1/\Delta f$. Das Instrument stimmen heisst also, diese Differenz auf $0$
zu bringen.
\end{merksatz}
